%% kapa: exam=1 round=34 year=2023 period=2 form=A subject=accounting subject_ko=회계학 seq=22 number=62
%% source: Q-Net 감정평가사 기출문제 (cst00309) artlSeq=5211789 file=raw/1008/5211789/2023년도 제34회 감정평가사 1차시험 2교시.pdf
%% conversion: kordoc 4.24.51 -> tools/gen1.py (auto); review=proofread-sonnet; flags=none
\begin{question}{62}
리스에 관한 설명으로 옳은 것을 모두 고른 것은?
\begin{framedtext}
ㄱ. 단기리스나 소액 기초자산 리스를 제외한 모든 리스에 대해서 리스이용자는 사용권자산과 리스부채를 인식해야 한다. \par
ㄴ. 리스이용자는 리스의 내재이자율을 쉽게 산정할 수 없는 경우에는 리스제공자의 증분차입이자율을 사용하여 리스료를 할인한다. \par
ㄷ. 리스이용자는 사용권자산이 손상되었는지를 판단하고 식별되는 손상차손을 회계처리하기 위하여 자산손상 기준서를 적용한다. \par
ㄹ. 투자부동산의 정의를 충족하는 사용권자산은 재무상태표에 투자부동산으로 표시한다. \par
\end{framedtext}
\begin{choices}
  \item ㄱ, ㄴ
  \item ㄱ, ㄷ
  \item ㄷ, ㄹ
  \item ㄱ, ㄷ, ㄹ
  \item ㄴ, ㄷ, ㄹ
\end{choices}
\end{question}
\answer{4}
